\clearpage
\section{정규성과 분리성의 결합}
\label{sec:galois-extensions}

\begin{learninggoal}
갈루아 확장을 정규성과 분리성의 결합으로 정의한다. 유한 갈루아 확장을 분리다항식의 분해체로 판정하고, 앞 장에서 만난 대표 확장들을 갈루아성의 관점에서 다시 분류한다.
\end{learninggoal}

앞 장까지 우리는 대수적 확장의 두 성질을 따로 연구했다.

\begin{itemize}
  \item \emph{분리성}은 최소다항식이 충분히 많은 서로 다른 근을 갖는다는 조건이다.
  \item \emph{정규성}은 한 켤레근이 확장체에 들어오면 나머지 켤레근도 모두 함께 들어온다는 조건이다.
\end{itemize}

갈루아 이론은 이 두 조건이 동시에 성립하여, 서로 다른 켤레근들이 모두 같은 체 안에서 자동동형으로 서로 이동하는 상황을 연구한다.

\begin{definition}[갈루아 확장과 갈루아군]
대수적 확장 $L/K$가 정규확장이면서 분리확장이면 $L/K$를 \emph{갈루아 확장}(Galois extension)이라 한다. 이때
\[
  \Gal(L/K):=\Aut_K(L)
\]
를 $L/K$의 \emph{갈루아군}(Galois group)이라 한다.
\end{definition}

정규성만으로는 자동동형들이 기준체를 완전히 드러내지 못할 수 있다. 비자명한 순수비분리 확장은 정규이지만 자동동형군이 항등원 하나뿐이다. 반대로 분리성만으로는 켤레근들이 확장체 밖으로 빠져나갈 수 있다. $\Q(\sqrt[3]{2})/\Q$가 그 예이다. 갈루아 확장에서는 두 결함이 모두 사라진다.

\begin{theorem}[유한 갈루아 확장의 분해체 판정]
유한 대수적 확장 $L/K$에 대해 다음은 동치이다.
\begin{enumerate}[label=(\roman*)]
  \item $L/K$는 갈루아 확장이다.
  \item $L$은 어떤 분리다항식 $f\in K[X]$의 분해체이다.
  \item $L$은 유한개의 분리다항식으로 이루어진 다항식족의 분해체이다.
\end{enumerate}
\end{theorem}

\begin{proof}
(i)$\Rightarrow$(ii)를 보자. 유한확장이므로
\[
  L=K(\alpha_1,\dots,\alpha_r)
\]
로 쓸 수 있다. 각 $\alpha_i$의 $K$ 위 최소다항식을 $m_i$라 하자. $L/K$가 분리이므로 모든 $m_i$는 분리다항식이고, 정규이므로 각 $m_i$는 $L$에서 완전히 분해된다. 중복되는 최소다항식을 한 번씩만 남기고 그 곱을
\[
  f=m_{i_1}\cdots m_{i_s}
\]
라 하자. 서로 다른 일계수 기약다항식들은 서로소이고 각각 분리이므로 $f$도 분리다항식이다. $f$의 모든 근은 $L$에 들어 있고, 그 근들 가운데 $\alpha_1,\dots,\alpha_r$가 있으므로 $f$의 분해체는 정확히 $L$이다.

(ii)$\Rightarrow$(iii)는 자명하다.

(iii)$\Rightarrow$(i)를 보자. 분해체는 정규확장이다. 또한 $L$은 분리다항식들의 근들로 생성되므로 분리확장이다. 따라서 $L/K$는 갈루아 확장이다.
\end{proof}

\begin{corollary}
분리다항식의 분해체는 유한 갈루아 확장이다. 특히 표수 $0$에서 다항식의 중근을 제거한 뒤 얻는 분해체는 갈루아 확장이다.
\end{corollary}

\begin{example}[이차확장]
$d\in\Q$가 제곱이 아니라고 하자. 확장
\[
  \Q(\sqrt d)/\Q
\]
는 $X^2-d$의 분해체이고 표수 $0$에서 분리이므로 갈루아 확장이다. 갈루아군은
\[
  \id,\qquad
  \sigma:\sqrt d\longmapsto-\sqrt d
\]
의 두 원소로 이루어진 순환군이다.
\end{example}

\begin{example}[분리이지만 비정규]
$\alpha=\sqrt[3]{2}$라 하자. $X^3-2$는 $\Q$ 위에서 분리이지만
\[
  \Q(\alpha)
\]
에는 다른 두 근 $\omega\alpha,\omega^2\alpha$가 들어 있지 않다. 따라서 $\Q(\alpha)/\Q$는 분리확장이지만 갈루아 확장은 아니다.
\end{example}

\begin{example}[정규이지만 비분리]
표수 $p>0$에서
\[
  \F_p(t)/\F_p(t^p)
\]
는 순수비분리이므로 정규확장이지만 분리확장이 아니다. 자동동형은 항등원 하나뿐이고, 확장차수는 $p$이다. 따라서 갈루아 확장이 아니다.
\end{example}

\begin{example}[유한체 확장]
모든 유한체 확장
\[
  \F_{q^n}/\F_q
\]
는 갈루아 확장이다. 실제로 $\F_{q^n}$은 분리다항식
\[
  X^{q^n}-X
\]
의 $\F_q$ 위 분해체이다. 그 갈루아군은 상대 프로베니우스
\[
  \Frob_q:a\longmapsto a^q
\]
가 생성하는 위수 $n$의 순환군이다.
\end{example}

\begin{keyidea}
유한 갈루아 확장은 \emph{분리된 모든 근을 한 체 안에 모아 놓은 확장}이다. 따라서 갈루아군은 그 근들을 서로 바꾸되 기준체의 원소는 움직이지 않는 대칭들의 군이다.
\end{keyidea}

\begin{checkpoint}
\begin{enumerate}[label=(\alph*)]
  \item $\Q(i)/\Q$가 갈루아 확장임을 분해체 판정으로 설명하라.
  \item $\Q(\sqrt[4]{2})/\Q$가 분리이지만 갈루아가 아닌 이유를 적어라.
  \item $\F_{q^6}/\F_q$의 갈루아군의 위수와 한 생성원을 구하라.
\end{enumerate}
\end{checkpoint}
